% !TeX TS-program = xelatex
% 虚构演示简历 — 仅用于 hellojobs 测试数据
\documentclass{resume-photo}
\usepackage{xeCJK}
\setCJKmainfont{Noto Serif CJK SC}
\usepackage{fontspec}
\usepackage{enumitem}
\setlist[itemize]{itemsep=0pt, topsep=1pt, parsep=0pt, partopsep=0pt}

\ResumeName{何森}

\begin{document}

\ResumeContacts{
  138-0015-0015,%
  \ResumeUrl{mailto:hesen@bit.edu.cn}{hesen@bit.edu.cn},%
  \textnormal{北京理工大学 | 软件工程 · 硕士 | 2023—2026}%
}

\ResumeTitle

\noindent 区块链与智能合约，熟悉 Solidity、Hardhat 与 EVM 安全最佳实践；参与联盟链存证模块链码审计辅助。注重可观测性与文档沉淀，习惯在评审阶段拆解风险；参与线上复盘与跨团队联调，英语 CET-6，可阅读官方文档与 RFC（演示数据）。

\section{教育背景}
\ResumeItem
{北京理工大学}
[\textnormal{计算机学院 | 软件工程 | 硕士}]
[2023.09—2026.06(预计)]

\begin{itemize}
  \item 分布式系统与可信计算交叉方向。
\end{itemize}


\ResumeItem
{企业开放日与实训}
[\textnormal{短期实训营（演示）}]
[2023—2024]

\begin{itemize}
  \item 参加头部互联网公司 2 周封闭实训，完成从需求到上线的完整小组项目。
  \item 在实训中负责接口设计与联调，小组答辩成绩排名前 20\%。
  \item 获得实训营「优秀学员」与内推绿色通道（测试数据，勿当真）。
  \item 沉淀实训笔记 1.5 万字，含架构图与排错记录。
\end{itemize}



\section{专业技能}
\begin{itemize}
  \item 熟悉 Solidity、OpenZeppelin、Hardhat/Foundry 测试。
  \item 了解 Hyperledger Fabric 链码生命周期与 CouchDB 状态模型。
  \item 掌握常见重入、整数溢出（历史）与访问控制漏洞排查。
  \item 熟悉 Linux 日常运维与 Shell，具备日志检索与 CPU/内存突增初步定位能力。
  \item 掌握 Git / CI 与 Code Review，了解 Docker 与 Kubernetes 基本编排与滚动发布。
  \item 了解 OAuth2 / JWT、接口幂等与 REST 规范，具备单元测试与集成测试习惯（JUnit/pytest 等）。
  \item 能阅读英文技术文档，具备跨团队沟通与需求澄清能力（测试简历）。
\end{itemize}

\section{实习经历}
\ResumeItem
{某金融科技公司}
[\textnormal{区块链研发实习生}]
[2024.07—2024.12]

\begin{itemize}
  \item \textbf{职责}: 存证合约升级与 Gas 优化。
  \item \textbf{成果}: 单笔存证 Gas 消耗下降约 18\%。
\end{itemize}


\ResumeItem
{网易 | 杭州研究院}
[\textnormal{中间件方向实习生}]
[2023.07—2024.01]

\begin{itemize}
  \item \textbf{消息}: 参与 MQ 控制台 Topic 治理与消息轨迹查询性能优化。
  \item \textbf{存储}: 调研并输出对象存储冷热分层方案对比报告 25 页。
  \item \textbf{演练}: 参与双活切换演练，验证 RPO/RTO 指标达标。
  \item \textbf{开源}: 向社区组件提交文档 PR 2 个并被合并。
  \item \textbf{软技能}: 跨 3 团队协调排期，保证里程碑按时交付。
\end{itemize}



\section{项目经历}
\ResumeItem
{NFT 票务 Demo}
[\textnormal{课程项目}]
[2024.02—2024.06]

\begin{itemize}
  \item ERC-721 + 链下元数据 IPFS Pinning，支持二级市场版税 5\%。
\end{itemize}


\ResumeItem
{实时风控规则引擎}
[\textnormal{Drools / 自研 DSL（演示）}]
[2024.01—2024.05]

\begin{itemize}
  \item \textbf{需求}: 支持运营侧快速上线规则，T+0 生效且可审计。
  \item \textbf{架构}: 规则编译为字节码缓存，热加载版本号隔离。
  \item \textbf{指标}: 规则平均评估耗时 <2ms，峰值 QPS 12k 稳定。
  \item \textbf{治理}: 规则变更双人复核 + 影子流量对比。
\end{itemize}



\section{个人荣誉}
\begin{itemize}
  \item 北京理工大学 学业奖学金
  \item CCF 区块链专委会 暑期学校 结业
  \item 全国计算机等级考试四级（网络工程师）（演示）
  \item 校级「优秀学生干部 / 优秀团员」或技术社团骨干（综合测评前 15\%）
  \item 开源贡献：向知名项目提交被合并的 PR（文档/测试/feature）（演示）
\end{itemize}

\end{document}
